% GENERATED FILE. Edit canonical lessons, then run npm run generate:books.
% canonical-source-hash: fnv1a64:d87585de7fcbd51f
% canonical-lessons: PA-C265-ure, PA-C265-pare, PA-C265-vichkar, PA-C265-vat, PA-C265-koh

\chapter{This Side and That}
\label{ch:pa-a2-this-side-and-that}
\hlchaptermodality{265}

\begin{chapteropening}
\textbf{By the end of this chapter:} \emph{I can say how near or far a place is.}

Complete the last lesson of chapter 265: I can say how near or far a place is.
\end{chapteropening}

\section[\texorpdfstring{\pa{ਉਰੇ} (ure)}{ਉਰੇ (ure)}]{\pa{ਉਰੇ} (ure) — on this side; nearer}

\label{lesson:PA-C265-ure}



\begin{quote}
Before the new one: say the Punjabi for a pole, then the Punjabi for a corner.
\end{quote}

\subsection*{You'll want to know: \pa{ਉਰੇ}}
\textbf{\pa{ਉਰੇ}} — \emph{ure} — \textquotedblleft{}on this side; nearer\textquotedblright{}.

\begin{quote}
\pa{ਬੱਸ} \pa{ਪੁਲ} \pa{ਤੋਂ} \pa{ਉਰੇ} \pa{ਰੁਕਦੀ} \pa{ਹੈ।}
\end{quote}

\begin{quote}
\emph{bass pul ton ure rukdī hai.}
\end{quote}

\begin{quote}
The bus stops on this side of the bridge.
\end{quote}

\begin{cousinweb}
\textbf{\pa{ਉਰੇ}} (\emph{ure}) is inherited from Sanskrit \emph{avara}, \textquotedblleft{}lower, hinder, nearer\textquotedblright{}; CDIAL lists Punjabi \emph{ure}, \textquotedblleft{}on this side\textquotedblright{}, under it (CDIAL 805).
\end{cousinweb}

\begin{grammarlens}[title={What you've built}]
One more word for this chapter.
\end{grammarlens}

\subsection*{Your turn}
\begin{itemize}
\raggedright
  \item \emph{Say it:} \emph{ure}
  \item \emph{Say it:} \emph{ure}, once more
  \item \emph{Say it:} \emph{bass pul ton ure rukdī hai.} — The bus stops on this side of the bridge.
  \item \emph{Recall:} say the Punjabi for a pole, then the Punjabi for a corner, then say \emph{ure} again
\end{itemize}

\begin{tcolorbox}[breakable,colback=teal!4,colframe=teal!35!black,title={Before you move on}]
What does \pa{ਉਰੇ} mean? (\textquotedblleft{}On this side; nearer\textquotedblright{}.) Say it once more.
\end{tcolorbox}
\section[\texorpdfstring{\pa{ਪਰੇ} (pare)}{ਪਰੇ (pare)}]{\pa{ਪਰੇ} (pare) — farther off; beyond}

\label{lesson:PA-C265-pare}



\begin{quote}
Before the new one: say the Punjabi for a corner, then the Punjabi for on this side.
\end{quote}

\subsection*{You'll want to know: \pa{ਪਰੇ}}
\textbf{\pa{ਪਰੇ}} — \emph{pare} — \textquotedblleft{}farther off; beyond\textquotedblright{}.

\begin{quote}
\pa{ਸਟੇਸ਼ਨ} \pa{ਪਿੰਡ} \pa{ਤੋਂ} \pa{ਬਹੁਤ} \pa{ਪਰੇ} \pa{ਹੈ।}
\end{quote}

\begin{quote}
\emph{saṭeshan piṇḍ ton bahut pare hai.}
\end{quote}

\begin{quote}
The station is a long way beyond the village.
\end{quote}

\begin{cousinweb}
\textbf{\pa{ਪਰੇ}} (\emph{pare}) is inherited from Sanskrit \emph{para}, \textquotedblleft{}distant, further, other\textquotedblright{}; CDIAL lists Punjabi \emph{pare}, \textquotedblleft{}beyond, further\textquotedblright{}, under it (CDIAL 7793).
\end{cousinweb}

\begin{grammarlens}[title={What you've built}]
One more word for this chapter.
\end{grammarlens}

\subsection*{Your turn}
\begin{itemize}
\raggedright
  \item \emph{Say it:} \emph{pare}
  \item \emph{Say it:} \emph{pare}, once more
  \item \emph{Say it:} \emph{saṭeshan piṇḍ ton bahut pare hai.} — The station is a long way beyond the village.
  \item \emph{Recall:} say the Punjabi for a corner, then the Punjabi for on this side, then say \emph{pare} again
  \item \emph{Write it:} copy \textbf{\pa{ਪਰੇ}} once, checking the whole word
\end{itemize}

\begin{tcolorbox}[breakable,colback=teal!4,colframe=teal!35!black,title={Before you move on}]
What does \pa{ਪਰੇ} mean? (\textquotedblleft{}Farther off; beyond\textquotedblright{}.) Say it once more.
\end{tcolorbox}
\section[\texorpdfstring{\pa{ਵਿਚਕਾਰ} (vichkār)}{ਵਿਚਕਾਰ (vichkār)}]{\pa{ਵਿਚਕਾਰ} (vichkār) — in the middle; between}

\label{lesson:PA-C265-vichkar}



\begin{quote}
Before the new one: say the Punjabi for on this side, then the Punjabi for farther off.
\end{quote}

\subsection*{You'll want to know: \pa{ਵਿਚਕਾਰ}}
\textbf{\pa{ਵਿਚਕਾਰ}} — \emph{vichkār} — \textquotedblleft{}in the middle; between\textquotedblright{}.

\begin{quote}
\pa{ਬੈਂਕ} \pa{ਸਕੂਲ} \pa{ਅਤੇ} \pa{ਦਫ਼ਤਰ} \pa{ਦੇ} \pa{ਵਿਚਕਾਰ} \pa{ਹੈ।}
\end{quote}

\begin{quote}
\emph{baiṅk sakūl ate daftar de vichkār hai.}
\end{quote}

\begin{quote}
The bank is between the school and the office.
\end{quote}

\begin{cousinweb}
\textbf{\pa{ਵਿਚਕਾਰ}} (\emph{vichkār}) is built on \emph{vicc}, \textquotedblleft{}middle\textquotedblright{}, from a reconstructed \emph{*vīcya}, \textquotedblleft{}middle\textquotedblright{}; CDIAL lists Punjabi \emph{vicc} under it (CDIAL 12042). How the ending \emph{-kār} joined it is not settled.
\end{cousinweb}

\begin{grammarlens}[title={What you've built}]
One more word for this chapter.
\end{grammarlens}

\subsection*{Your turn}
\begin{itemize}
\raggedright
  \item \emph{Say it:} \emph{vichkār}
  \item \emph{Say it:} \emph{vichkār}, once more
  \item \emph{Say it:} \emph{baiṅk sakūl ate daftar de vichkār hai.} — The bank is between the school and the office.
  \item \emph{Recall:} say the Punjabi for on this side, then the Punjabi for farther off, then say \emph{vichkār} again
  \item \emph{Write it:} copy \textbf{\pa{ਵਿਚਕਾਰ}} once, checking the whole word
\end{itemize}

\begin{tcolorbox}[breakable,colback=teal!4,colframe=teal!35!black,title={Before you move on}]
What does \pa{ਵਿਚਕਾਰ} mean? (\textquotedblleft{}In the middle; between\textquotedblright{}.) Say it once more.
\end{tcolorbox}
\section[\texorpdfstring{\pa{ਵਾਟ} (vāṭ)}{ਵਾਟ (vāṭ)}]{\pa{ਵਾਟ} (vāṭ) — the distance to travel; the way}

\label{lesson:PA-C265-vat}



\begin{quote}
Before the new one: say the Punjabi for farther off, then the Punjabi for in the middle.
\end{quote}

\subsection*{You'll want to know: \pa{ਵਾਟ}}
\textbf{\pa{ਵਾਟ}} — \emph{vāṭ} — \textquotedblleft{}the distance to travel; the way\textquotedblright{}.

\begin{quote}
\pa{ਸ਼ਹਿਰ} \pa{ਦੀ} \pa{ਵਾਟ} \pa{ਲੰਮੀ} \pa{ਹੈ।}
\end{quote}

\begin{quote}
\emph{shahir dī vāṭ lammī hai.}
\end{quote}

\begin{quote}
The way to the town is long.
\end{quote}

A feminine noun.

\begin{cousinweb}
\textbf{\pa{ਵਾਟ}} (\emph{vāṭ}) is inherited from Sanskrit \emph{vartman}, \textquotedblleft{}track of a wheel, path\textquotedblright{}; CDIAL lists Punjabi \emph{vāṭ}, \textquotedblleft{}path, road\textquotedblright{}, under it (CDIAL 11366).
\end{cousinweb}

\begin{grammarlens}[title={What you've built}]
One more word for this chapter.
\end{grammarlens}

\subsection*{Your turn}
\begin{itemize}
\raggedright
  \item \emph{Say it:} \emph{vāṭ}
  \item \emph{Say it:} \emph{vāṭ}, once more
  \item \emph{Say it:} \emph{shahir dī vāṭ lammī hai.} — The way to the town is long.
  \item \emph{Recall:} say the Punjabi for farther off, then the Punjabi for in the middle, then say \emph{vāṭ} again
  \item \emph{Write it:} copy \textbf{\pa{ਵਾਟ}} once, checking the whole word
\end{itemize}

\begin{tcolorbox}[breakable,colback=teal!4,colframe=teal!35!black,title={Before you move on}]
What does \pa{ਵਾਟ} mean? (\textquotedblleft{}The distance to travel; the way\textquotedblright{}.) Say it once more.
\end{tcolorbox}
\section[\texorpdfstring{\pa{ਕੋਹ} (koh)}{ਕੋਹ (koh)}]{\pa{ਕੋਹ} (koh) — a kos, about two miles}

\label{lesson:PA-C265-koh}



\begin{quote}
Before the new one: say the Punjabi for in the middle, then the Punjabi for the distance to travel.
\end{quote}

\subsection*{You'll want to know: \pa{ਕੋਹ}}
\textbf{\pa{ਕੋਹ}} — \emph{koh} — \textquotedblleft{}a kos, about two miles\textquotedblright{}.

\begin{quote}
\pa{ਮੇਰਾ} \pa{ਪਿੰਡ} \pa{ਸ਼ਹਿਰ} \pa{ਤੋਂ} \pa{ਦੋ} \pa{ਕੋਹ} \pa{ਦੂਰ} \pa{ਹੈ।}
\end{quote}

\begin{quote}
\emph{merā piṇḍ shahir ton do koh dūr hai.}
\end{quote}

\begin{quote}
My village is two kos from the town.
\end{quote}

A masculine noun.

\begin{cousinweb}
\textbf{\pa{ਕੋਹ}} (\emph{koh}) is inherited from Sanskrit \emph{krośa}, \textquotedblleft{}a shout; a measure of distance, as far as a voice carries\textquotedblright{}; CDIAL lists Punjabi \emph{koh}, \textquotedblleft{}a distance of about two miles\textquotedblright{}, under it (CDIAL 3611).
\end{cousinweb}

\begin{grammarlens}[title={What you've built}]
One more word for this chapter.
\end{grammarlens}

\subsection*{Your turn}
\begin{itemize}
\raggedright
  \item \emph{Say it:} \emph{koh}
  \item \emph{Say it:} \emph{koh}, once more
  \item \emph{Say it:} \emph{merā piṇḍ shahir ton do koh dūr hai.} — My village is two kos from the town.
  \item \emph{Recall:} say the Punjabi for in the middle, then the Punjabi for the distance to travel, then say \emph{koh} again
  \item \emph{Write it:} copy \textbf{\pa{ਕੋਹ}} once, checking the whole word
\end{itemize}

\begin{tcolorbox}[breakable,colback=teal!4,colframe=teal!35!black,title={Before you move on}]
What does \pa{ਕੋਹ} mean? (\textquotedblleft{}A kos, about two miles\textquotedblright{}.) Say it once more.
\end{tcolorbox}
